The study was motivated by the need to make prostate-cancer AI reliability clinically
interrogable: which quantitative axes are represented, which remain stable across cohort and
technical changes, and which are used by a downstream decision. The comparative hierarchy
answers the first two questions across six non-interchangeable axes. Grade/ISUP and tumor
phenotype/content were the strongest transported representation coordinates; PTEN loss and AR
activity were conditional; SPOP mutation was unsupported in the frozen design; and recurrence
was endpoint- and reference-sensitive. The third question advanced only to a scoped internal
ISUP functional-sensitivity example and did not establish indispensable or externally transported
head use.

The evidence chain supporting that conclusion is deliberately asymmetric. Unchanged-probe
transport makes morphology-linked axes candidate cross-cohort audit anchors; conditional and
site analyses show that PTEN and AR decodability does not establish stable independent
explanation; an available genomic comparison rejects SPOP as a qualified frozen-design
coordinate; and endpoint contrasts prevent recurrence association from being treated as an
endpoint-invariant prognostic marker. Only the ISUP erasure experiment connects a recoverable
direction to sensitivity of a locked downstream judgment, while its refit results prevent a claim
of indispensable use.

This qualified map provides contestable coordinates for model audit. Transported coordinates
can anchor clinician review of agreement and disagreement, conditional or unsupported axes show
where explanations should be withheld, and evidence gaps prioritize external or functional
validation. Unavailable labels are not negative results. These uses do not demonstrate
pathologist agreement, calibrated reliance, clinical
benefit, or complete explanation of AI representations or decisions.

The map does not itself demonstrate a new residual marker; it defines a disciplined starting
point for potential new-biomarker discovery. After jointly accounting for known
clinical--pathology targets and technical confounders, residual decision-linked signals must
recur across models and independent cohorts, contribute to a locked prediction head, become
reproducible human-measurable features, and receive biological or clinical validation before
being considered candidate biomarkers. External functional validation, clinician studies, and
residual-marker discovery therefore remain distinct follow-up studies.
